\begin{algorithm}[H]
\caption{Mixture Unscented Transform with Moment-Preserving Axial Splitting (M-UT)}
\label{alg:mut_moment_split}
\begin{algorithmic}[1]
\State \textbf{Inputs:} initial mean $\mu_0\in\mathbb{R}^d$, covariance $P_0\in\mathbb{R}^{d\times d}$; component count $N\ge 1$; split rank \(r\in\{1,2\}\); variance split $\alpha\in[0,1]$; horizon $\tau$; dynamics map $F_{t\to t+\tau}$; UT weights $(w^\mu,w^P)$
\State \textbf{Outputs:} propagated component moments $\{\mu_i^+,P_i^+\}_{i=1}^N$, mixture weights $W_i=\tfrac{1}{N}$, and recombined moments $(\mu^+,P^+)$
\Statex

\State \textbf{Eigenstructure and subspace split}
\State Compute $P_0 = U\,\mathrm{diag}(\lambda)\,U^\top$ with $\lambda_1\ge\dots\ge\lambda_d$
\State Set $r \gets \texttt{split\_rank}$; $U_r \gets U_{(:,1\!:\!r)}$, $\Lambda_r \gets \mathrm{diag}(\lambda_{1:r})$
\State $P_{\parallel} \gets U_r \Lambda_r U_r^\top$;\quad $P_{\perp} \gets P_0 - P_{\parallel}$
\State $\Sigma \gets (1-\alpha^2)\,P_{\parallel} + P_{\perp}$ \Comment{within-component covariance (common to all components)}

\Statex
\State \textbf{Isotropic point set in principal subspace}
\If{$r=0$ \textbf{or} $N=1$}
  \State $Z \gets \mathbf{0}_{0\times N}$
\ElsIf{$N=2r+1$}
  \State Axial layout: $Z\in\mathbb{R}^{r\times N}$ with columns $\{0,\pm c\,e_1,\ldots,\pm c\,e_r\}$, $c=\sqrt{N/2}$
\ElsIf{$N=r+1$}
  \State Regular simplex: build $Z$ from an orthonormal basis of the subspace orthogonal to $\mathbf{1}_N$ so that $\tfrac{1}{N}ZZ^\top=I_r$
\Else
  \State Random-whitened: draw $G\sim\mathcal{N}(0,I_{r\times N})$, center columns, then set $Z\gets C^{-1/2}G$ where $C=\tfrac{1}{N}GG^\top$
\EndIf
\State Ensure $\tfrac{1}{N}ZZ^\top=I_r$ (by construction in each case)

\Statex
\State \textbf{Component means, covariances, and weights}
\State $L_r \gets U_r\,\Lambda_r^{1/2}$
\State $\Delta \gets \alpha\,L_r Z\in\mathbb{R}^{d\times N}$
\For{$i=1$ \textbf{to} $N$}
  \State $\mu_i \gets \mu_0 + \Delta_{(:,i)}$
  \State $P_i \gets \Sigma$
  \State $W_i \gets 1/N$
\EndFor
\If{verification is enabled}
  \State Check $\sum_i W_i\mu_i = \mu_0$ and $\sum_i W_i\!\big(P_i+(\mu_i-\mu_0)(\mu_i-\mu_0)^\top\big)=P_0$ (raise on failure)
\EndIf

\Statex
\State \textbf{Per-component UT propagation}
\For{$i=1$ \textbf{to} $N$}
  \State Generate $(2d+1)$ $\sigma$-points $\{\chi_{i,c}\}$ for $\mathcal{N}(\mu_i,P_i)$ with weights $(w^\mu,w^P)$
  \State Propagate: $\chi_{i,c}^+ \gets F_{t\to t+\tau}(\chi_{i,c})$
  \State Rebuild: $\mu_i^+ \gets \sum_{c} w^\mu \chi_{i,c}^+$;\quad
         $P_i^+ \gets \sum_{c} w^P(\chi_{i,c}^+-\mu_i^+)(\chi_{i,c}^+-\mu_i^+)^\top$
\EndFor

\Statex
\State \textbf{Mixture recombination (exact first two moments)}
\State $\mu^+ \gets \sum_{i=1}^N W_i\,\mu_i^+$
\State $P^+ \gets \sum_{i=1}^N W_i\!\Big(P_i^+ + (\mu_i^+-\mu^+)(\mu_i^+-\mu^+)^\top\Big)$

\State \textbf{Return} $\{\mu_i^+,P_i^+\}_{i=1}^N$, $\{W_i\}_{i=1}^N$, and $(\mu^+,P^+)$
\end{algorithmic}
\end{algorithm}


%       File: VTthesis_template.tex
%     Created: Thu Mar 24 11:00 AM 2016 EDT
%     Last Change: August 15, 2024
%     Author: Alan M. Lattimer, VT
%	  With modifications by Carrie Cross, Robert Browder, and LianTze Lim. 
%
% This template is designed to operate with XeLaTeX.
%
% All elements in the Title, Abstract, and Keywords MUST be formatted as text and NOT as math.
%
%Further instructions for using this template are embedded in the document. Additionally, there are comments at the end of the file that give suggestions on writing your thesis.  
%
%In addition to the standard formatting options, the following options are defined for the VTthesis class: proposal, prelim, doublespace, draft. 



%****************************************************************************
% Below are some general suggestions for writing your dissertation:
%
% 1. Label everything with a meaningful prefix so that you
%    can refer back to sections, tables, figures, equations, etc.
%    Usage \label{<prefix>:<label_name>} where some suggested
%    prefixes are:
%			ch: Chapter
%     		se: Section
%     		ss: Subsection
%     		sss: Sub-subsection
%			app: Appendix
%     		ase: Appendix section
%     		tab: Tables
%     		fig: Figures
%     		sfig: Sub-figures
%     		eq: Equations
%
% 2. The VTthesis class provides for natbib citations. You should upload
%	 one or more *.bib bibtex files. Suppose you have two bib files: some_refs.bib and 
%    other_refs.bib.  Then your bibliography line to include them
%    will be:
%      \bibliography{some_refs, other_refs}
%    where multiple files are separated by commas. In the body of 
%    your work, you can cite your references using natbib citations.
%    Examples:
%      Citation                     Output
%      =======================================================
%      \cite{doe_title_2016}        [18]
%      \citet{doe_title_2016}       Doe et al. [18]
%      \citet*{doe_title_2016}      Doe, Jones, and Smith [18]
%
%    For a complete list of options, see
%      https://www.ctan.org/pkg/natbib?lang=en
%
% 3. Here is a sample table. Notice that the caption is centered at the top. Also
%    notice that we use booktabs formatting. You should not use vertical lines
%    in your tables.
% 
%				\begin{table}[htb]
%					\centering
%					\caption{Approximate computation times in hh:mm:ss for full order 						versus reduced order models.}
%					\begin{tabular}{ccc}
%						\toprule
%						& \multicolumn{2}{c}{Computation Time}\\
%						\cmidrule(r){2-3}
%						$\overline{U}_{in}$ m/s & Full Model & ROM \\
%						\midrule
%						0.90 & 2:00:00 & 2:08:00\\
%						0.88 & 2:00:00 & 0:00:03\\
%						0.92 & 2:00:00 & 0:00:03\\
%						\midrule
%						Total & 6:00:00 & 2:08:06\\
%						\bottomrule
%					\end{tabular}
%					\label{tab:time_rom}
%				\end{table}
% 
% 4. Below are some sample figures. Notice the caption is centered below the
%    figure.
%    a. Single centered figure:
%					\begin{figure}[htb]
%						\centering
%						\includegraphics[scale=0.5]{my_figure.eps}
%						\caption{Average outlet velocity magnitude given an average  
%				        input velocity magnitude of 0.88 m/s.} 
%						\label{fig:output_rom}
%					\end{figure}
%    b. Two by two grid of figures with subcaptions
%					\begin{figure}[htb]
%						\centering
%						\begin{subfigure}[h]{0.45\textwidth}
%							\centering
%							\includegraphics[scale=0.4]{figure_1_1.eps}
%							\caption{Subcaption number one}
%							\label{sfig:first_subfig}
%						\end{subfigure}
%						\begin{subfigure}[h]{0.45\textwidth}
%							\centering
%							\includegraphics[scale=0.4]{figure_1_2.png}
%							\caption{Subcaption number two}
%							\label{sfig:second_subfig}
%						\end{subfigure}
%
%						\begin{subfigure}[h]{0.45\textwidth}
%							\centering
%							\includegraphics[scale=0.4]{figure_2_1.pdf}
%							\caption{Subcaption number three}
%							\label{sfig:third_subfig}
%						\end{subfigure}
%						\begin{subfigure}[h]{0.45\textwidth}
%							\centering
%							\includegraphics[scale=0.4]{figure_2_2.eps}
%							\caption{Subcaption number four}
%							\label{sfig:fourth_subfig}
%						\end{subfigure}
%						\caption{Here is my main caption describing the relationship between the 4 subimages}
%						\label{fig:main_figure}
%					\end{figure}
%
%============================================================================
%
% The following is a list of definitions and packages provided by VTthesis:
%
% A. The following packages are provided by the VTthesis class:
%      amsmath, amsthm, amssymb, enumerate, natbib, hyperref, graphicx, 
%      tikz (with shapes and arrows libraries), caption, subcaption,
%      listings, verbatim
%
% B. The following theorem environments are defined by VTthesis:
%      theorem, proposition, lemma, corollary, conjecture
% 
% C. The following definition environments are defined by VTthesis:
%      definition, example, remark, algorithm
%
%============================================================================
%
%  I hope this template file and the VTthesis class will keep you from having 
%  to worry about the formatting and allow you to focus on the actual writing.
%  Good luck, and happy writing.
%    Alan Lattimer, VT, 2016
%
%****************************************************************************
